\subsection{Completion of the level complexes}

Fix a finite-index pure subgroup $G\subset\Gamma_{g,n}$ so that the simplicial action has no inversions; taking a barycentric subdivision also removes inversions. For a simplicial $G$-complex $K$ with finitely many cell orbits, its \emph{full level completion} is the simplicial profinite set
\[
 K^{\mathrm{lev}}=\varprojlim_{H\leq_fG}K/H.
\]
The procongruence version uses only congruence subgroups $H$. The map between them is induced by restriction of the inverse system. These definitions fix the exact completions used below.

\begin{result}{D.11}{Proved}{Profinite simple connectivity at level two and above}
For every stable type and every $k\geq2$, the full level completion of $P_k$ has trivial profinite fundamental group. The completed $P_1$ is connected. For the full curve complex, the same simple-connectivity conclusion holds for $d\geq3$. Completing the discrete presentation from the complex of groups on $P_2$ gives the full profinite group $\widehat\Gamma_{g,n}$ as a profinite presentation, with generators furnished by strata of modular dimension at most one and additional gluing relations furnished in dimension two.
\end{result}
\begin{proof}
D.5 and the cone convention show that $P_k$ is simply connected for $k\geq2$. For any finite-index $H\subset G$, Armstrong's quotient theorem gives
\[
 \pi_1(P_k/H)=H/H_*,
\]
where $H_*$ is the normal subgroup generated by elements with fixed points. The theorem's hypotheses hold: the complex is simply connected, the action is simplicial and without inversions. The same theorem and derivation are opened in Kent \S7, Theorem~11 \cite{Kent2016}. Thus $\widehat H$ surjects onto $\widehat{\pi_1(P_k/H)}$. Inverse limits of these surjections are surjective by compactness. But $\varprojlim_{H\leq_fG}\widehat H=1$: each $\widehat H$ identifies with the closure of $H$ in $\widehat G$, and their intersection is the identity by the defining finite-quotient topology. The profinite fundamental group commutes with this inverse limit of finite simplicial sets, in the sense specified in Kent \S5. Hence the desired group is trivial. The proof for the curve complex uses its simple connectivity at $d\geq3$. Connectedness of completed $P_1$ follows from connectedness of all its finite-level quotients and the corresponding inverse-limit statement for $\pi_0$.

For the presentation assertion, choose a maximal tree, stabilizer lifts, and transporter elements in the discrete complex of groups. Its ordinary fundamental group is $\Gamma_{g,n}$ because $P_2$ is simply connected. The homomorphisms from this presentation to any finite group are exactly the compatible homomorphisms from the local stabilizers and the edge transporters satisfying the edge and triangle relations. By the universal property of completion, homomorphisms from each local group to a finite group extend uniquely and continuously to its full completion. The same finite-group-valued universal property therefore identifies the resulting profinite presentation with $\widehat\Gamma_{g,n}$. This proves the presentation as a profinite colimit; it does not silently assert injectivity of every completed local-group map. Such injectivity is a separate comparison supplied where needed by Part~E.
\end{proof}

\begin{result}{D.12}{Refuted}{A printed low-complexity range in the profinite literature}
For the unmodified disjointness complex, the full level completion of $\mathcal C_{0,5}$ has a nontrivial continuous $\mathbb Z/2$-valued first cohomology class, and hence nontrivial profinite fundamental group. Thus the bound $3g-3+n\geq2$ printed in Kent's Theorem~12 cannot hold for the complex defined there. The argument in D.11 proves the corrected range $3g-3+n\geq3$. The genus-zero case of the analogous printed range in Theorem~13 likewise requires this correction.
\end{result}
\begin{proof}
Place the five labelled points on a circle. Let $c_{12},c_{23},c_{34},c_{45},c_{51}$ be boundaries of small regular neighborhoods of the five consecutive arcs. Curves associated with disjoint pairs are disjoint. Therefore
\[
 c_{12},c_{34},c_{51},c_{23},c_{45},c_{12}
\]
is a closed pentagon in the actual infinite curve complex. The quotient by the pure mapping class group maps a curve to the two-element side of its partition of the labels. It is the Petersen graph: edges join disjoint two-element subsets. The pentagon above maps to an embedded pentagon in that graph. On the Petersen graph take the $\mathbb Z/2$ cellular one-cochain equal to one on the edge $\{12,34\}$ and zero on all other edges. Every such cochain is a cocycle, since the graph has no two-cells, and its evaluation on this pentagon is one.

For every finite-index subgroup $H$ of the pure mapping class group, there are compatible maps
\[
 \mathcal C_{0,5}\longrightarrow\mathcal C_{0,5}/H
 \longrightarrow\mathcal C_{0,5}/\Gamma_{0,5}.
\]
The displayed closed pentagon supplies a closed lift in every intermediate graph. Thus the pulled-back cocycle remains nonzero at every level. It defines a nonzero continuous cocycle on the inverse-limit simplicial profinite graph: if it were a coboundary, a continuous zero-cochain with finite target would factor through one finite level, where the pentagon evaluation rules out a coboundary. This gives the asserted continuous character of the profinite fundamental group.

The finite Petersen graph, the lifted pair pattern, and the cocycle evaluation are checked in \texttt{code/d\_boundary\_checks.py}. No computation on a finite quotient is being used to prove that an arbitrary infinite loop survives: we supplied the actual closed lift. The source of the printed problem is visible in Kent's Theorem~10, which gives the correct Harer sphere dimensions but writes the complex dimension as $3g-3+n$ and infers simple connectivity one complexity too early. Armstrong's theorem requires that missing simple connectivity. The same sphere theorem gives $\mathcal C_{0,5}\simeq\bigvee_\infty S^1$ and D.11 supplies the corrected elementary argument.
\end{proof}

\begin{result}{D.13}{Known}{The proved procongruence range and the remaining scope}
For $n\geq2$, $2g-2+(n-1)>0$, and $3g-3+n\geq3$, the procongruence curve complex has trivial profinite fundamental group. This follows from Kent's Theorem~13 with its proof restricted to the actual simply connected range. The present proof does not establish the stronger higher-connectivity claims for arbitrary procongruence $P_k$, nor does it infer the congruence subgroup property from this one complex.
\end{result}
\begin{proof}
Put $S=S_{g,n-1}$ and $\dot S=S_{g,n}$. The first two inequalities make $S$ a smooth affine hyperbolic curve; the third makes $\mathcal C(\dot S)$ simply connected, which is the exact hypothesis required for the Armstrong sequence used in the proof. We spell out the comparison employed in Kent's \S8. Let $K$ be the congruence kernel, and let $\kappa$ range through finite-index congruence subgroups of the pure level subgroup. Armstrong's theorem, profinite right exactness, and inverse-limit exactness give
\[
 \varprojlim_\kappa\widehat{\kappa_*}\longrightarrow K
 \longrightarrow\Pi_1(\check{\mathcal C}(\dot S))\longrightarrow1.
\]
Choose a curve $\gamma$ bounding a pair of pants containing the added puncture and another puncture. Kent's Theorem~3 and Corollary~4 prove that the products $B_\kappa\kappa_\gamma$ are congruence and cofinal, where $B$ is the point-pushing subgroup and $B_\kappa=B\cap\kappa$. In the opened proof of Lemma~15, closure separation of the stabilizer implies
\[
 \bigcap_\kappa\overline{B_\kappa h_\gamma}
       =\overline{h_\gamma}
\]
for every fixed congruence subgroup $h$. Cofinality then places $K$ in the image of $\varprojlim\widehat{\kappa_\gamma}$, which lies in the image of $\varprojlim\widehat{\kappa_*}$. The displayed quotient is trivial. This is the exact step of the cited proof that transfers to our conventions. It uses the affine Birman sequence and the congruence cofinality proved in the source, not an assertion that all finite-index subgroups are congruence. D.12 explains why we have stated the necessary complexity restriction explicitly.
\end{proof}
